\section{Le télégraphe (1794--1880)}\label{sec:telegraph}
% =============================================================================

\subsection{Le télégraphe optique et électrique}\label{subsec:telegraph}

La transmission d'information à distance n'attendit pas l'électricité, mais les raisons pour lesquelles elle devint, au cours du dix-neuvième siècle, un enjeu économique croissant doivent être posées d'entrée. Entre 1840 et 1914, le commerce international change d'échelle. Le commerce extérieur de la France passe de 2,5~milliards de francs courants en 1847 à 15~milliards en 1913~; celui de l'Angleterre de 13 à 35~milliards entre 1870 et 1913, celui de l'Allemagne de 5 à 25~milliards \parencite{mitchell_historical_2007}\footnote{Valeurs exprimées en francs-or courants par \textcite{mitchell_historical_2007}. Les ordres de grandeur de la croissance sont comparables en valeur réelle.}. Des navires chargent du coton à New York et le déchargent à Liverpool, mais l'information sur les cours, les commandes et les stocks voyage au même rythme qu'eux, quand elle ne prend pas de retard. La poste, qui repose sur le transfert physique du document, ne peut qu'aligner la vitesse de l'information sur celle des personnes\footnote{\textcite{juhasz_spinning_2018} le formulent avec précision~: avant le télégraphe, l'information ne pouvait voyager qu'à la vitesse du mode de transport physique le plus rapide. \textcite{carey_communication_1989} fait de cette indissociabilité entre communication et transport le point de départ de son analyse~: le télégraphe est la première technologie à séparer les deux.}. \dimmark{D6}{?}%
Or la croissance du volume des échanges accroît le nombre de décisions~--- acheter, vendre, acheminer, assurer~--- qui dépendent d'une information distante~; chaque décision prise avec un retard de dix jours sur les cours réels est une décision prise dans l'obscurité. Le problème n'est pas que l'information soit lente~; c'est que le système qu'elle doit coordonner grossit plus vite qu'elle.

C'est dans cet écart que le télégraphe va s'insérer. Le principe de la transmission électrique d'un signal était connu depuis les expériences de Lesage à Genève (1774) et de Sömmering à Munich (1809), mais ces systèmes étaient des impasses~: celui de Lesage exigeait autant de fils que de lettres dans l'alphabet, celui de Sömmering reposait sur un bullage électrochimique trop lent pour un usage opérationnel \parencite{beauchamp2001telegraphy}. C'est la découverte de l'électromagnétisme par Oersted à Copenhague en 1820 qui ouvrit les deux trajectoires viables~: d'un côté le galvanomètre, dont Schilling à Saint-Pétersbourg puis Gauss et Weber à Göttingen tirèrent un télégraphe à aiguilles~; de l'autre l'électro-aimant, dont Sturgeon à Londres puis Henry à Albany firent un actionneur capable de commander un bras métallique à distance \parencite{kooij_telegraph_2015}\footnote{Le relais électromagnétique de Henry est le composant clé~: il permettait de régénérer le signal sur de longues distances, rendant le réseau possible. C'est le même principe~--- un circuit électrique qui en commande un autre~--- qui servira dans les réseaux de distribution électrique de la génération suivante.}. Les deux lignées aboutirent indépendamment au télégraphe électrique viable en 1837, la première avec Cooke et Wheatstone en Angleterre, la seconde avec Morse en Amérique. La capacité technique existait donc, mais elle ne suffit pas à expliquer le déploiement~: il fallait aussi une demande solvable, un acteur prêt à financer le réseau, et un cadre institutionnel qui en autorise ou en commande la construction\footnote{On distingue ici, par commodité analytique, la capacité technique (l'existence d'un artefact fonctionnel) des conditions de son déploiement effectif (demande, financement, cadre réglementaire). Cette distinction est un choix de cadrage~: la sociologie des techniques a montré que la capacité technique est elle-même socialement construite~--- un appareil existe quand un réseau d'acteurs le fait exister. Nous la maintenons parce qu'elle organise la comparaison entre cas, en rendant visible le contraste entre une capacité partagée et des trajectoires divergentes.}. Ces conditions se manifestèrent de manière hétérogène~: demande militaire en France (la Convention avait besoin de coordonner un territoire en guerre), demande ferroviaire en Angleterre (la sécurité des trains sur voie unique), demande commerciale aux États-Unis (l'arbitrage sur les cours du coton), commande étatique en Prusse (le suivi du Parlement de Francfort depuis Berlin). Dans chaque cas, la forme institutionnelle du réseau~--- monopole d'État, concession privée, initiative ferroviaire~--- refléta l'intérêt de l'acteur qui le finançait. La question est alors de savoir si ces conditions d'émergence, radicalement différentes, produisent des formes organisationnelles tout aussi différentes, ou si les propriétés du réseau lui-même finissent par dominer les conditions initiales. La réponse n'est pas donnée d'avance~; les cas qui suivent la documenteront. Ce que l'on peut déjà noter, c'est que les composantes de la demande n'ont pas la même temporalité~: la coordination militaire en temps de guerre, immédiate et étatique~; la coordination commerciale en temps de paix, progressive et privée~; la coordination ferroviaire, tardive et organisationnelle.

Que la communication rapide à distance fût possible, le réseau de sémaphores de Claude Chappe l'avait déjà montré en France. Le réseau fut déployé à partir de 1793 non par curiosité technique mais parce que la Convention nationale, en guerre contre les puissances européennes, avait un besoin urgent de coordonner les mouvements militaires sur un territoire immense~; il couvrait 534~stations en 1844 et plus de cinq mille kilomètres de lignes \parencite{holzmann_early_1995, field_french_1994}. La ligne Paris-Lille transmettait le premier symbole d'un message en neuf minutes~; la ligne Paris-Strasbourg mettait six heures pour une dépêche de 25~mots \parencite[p.~87]{holzmann_early_1995}. Mais le principe restait mécanique et optique --- des bras articulés observés à la longue-vue, reproduits de tour en tour, inutilisables la nuit et par temps de brouillard. En 1838, l'administration comptait près de mille stationnaires pour une cinquantaine d'inspecteurs \parencite{ducamp_telegraphe_1869, field_french_1994}. \dimmark{D3}{?}%
Le poste dominant n'était pas le capital matériel~--- les tours étaient des constructions modestes~--- mais le travail humain, irréductiblement attaché à chaque n\oe{}ud~: doubler la longueur du réseau doublait le nombre d'opérateurs.

C'est ce rapport entre capital et travail que le télégraphe électrique allait modifier~--- et l'observation la plus frappante, documentée indépendamment dans chaque contexte national, est qu'il le fit partout de la même manière. Le télégraphe électrique substituait du capital physique~--- le fil métallique, les poteaux, les isolateurs, les batteries~--- au travail humain des opérateurs de tour\footnote{Le fait de cette substitution est documenté indépendamment pour chaque pays par \textcite{holzmann_early_1995}, \textcite{kieve_electric_1973}, \textcite{hubbard_cooke_1965} et \textcite{headrick_invisible_1991}. Aucun réseau télégraphique électrique ne conserva le principe des relais humains visuels.}. Il fonctionnait la nuit. Il fonctionnait par temps de brouillard. Sa disponibilité ne dépendait plus du climat mais de l'intégrité matérielle de la ligne. En contrepartie, il exigeait un investissement initial considérable en infrastructure linéaire~: les séries du \emph{National Bureau of Economic Research}\footnote{\emph{NBER} --- \emph{National Bureau of Economic Research}, organisme américain de recherche économique fondé en 1920, principal producteur de séries statistiques longues pour l'économie des États-Unis.} compilées par \textcite{gordon_capital_nineteenth} montrent que le réseau américain passa de 12\,000~miles de fil en 1850 à plus d'un million en 1900. La compétition entre les deux systèmes rivaux~--- le télégraphe à aiguilles de Cooke et Wheatstone\footnote{William Fothergill Cooke (1806--1879) et Charles Wheatstone (1802--1875), inventeurs britanniques du télégraphe électrique à aiguilles, breveté en 1837. Leur système utilisait initialement cinq aiguilles aimantées dont les déviations indiquaient les lettres, puis fut simplifié à deux aiguilles.}, qui nécessitait cinq fils puis deux \parencite{hubbard_cooke_1965}, et le système de Morse\footnote{Samuel Finley Breese Morse (1791--1872), peintre américain sans formation scientifique, conceptualisa le télégraphe en octobre~1832, sur le paquebot \emph{Sully} qui le ramenait de trois ans en Europe. Pendant les six semaines de traversée, le Dr.~Charles T.~Jackson de Boston lui montra un électro-aimant~; Morse esquissa son système dans son carnet de croquis. En débarquant à New York, il dit au capitaine Pell~: \og\emph{should you ever hear of the telegraph one of these days as the wonder of the world, remember the discovery was made on the good ship Sully}\fg{} \parencite[p.~251--257]{prime_life_1875}. Il passa cinq ans de misère à construire un prototype avec l'aide du mécanicien Alfred Vail (1807--1859), qui améliora considérablement le matériel et contribua à l'élaboration du code. Le système fut breveté en 1840~; la première dépêche publique fut transmise le 24~mai 1844 sur la ligne Washington-Baltimore~: \og\emph{What hath God wrought}\fg{} \parencite[p.~17]{bertho_histoire_1984}. Morse n'était pas électricien~; comme Cooke (ancien militaire), comme Bell (enseignant pour sourds), comme Marconi (autodidacte), il pensait en termes de fonction~--- communiquer à distance~--- pas en termes de substrat électrique \parencite{kooij_telegraph_2015}.}, qui n'en requérait qu'un~--- illustre la manière dont la structure de coûts orientait les choix techniques, sans les déterminer seule~: le code Morse économisait du cuivre, mais son succès tenait aussi à la qualité du matériel amélioré par Vail, au financement fédéral de la première ligne, et à la familiarité des opérateurs avec le système une fois installé~; d'autres trajectoires étaient possibles, et le télégraphe à aiguilles resta en usage en Grande-Bretagne jusqu'au \textsc{xx}\ieme{} siècle. Les opérateurs Morse devinrent, selon l'expression de \textcite{gabler_american_1988}, les premiers \og travailleurs de l'information\fg{} au sens propre~--- là où Wheatstone avait minimisé le besoin de compétence cognitive au prix d'un capital physique accru, \dimmark{D3/D4}{*}%
Morse fit l'inverse.

\dimmark{D2/D3}{*}%
La substitution du capital au travail dans la transmission d'information est documentée dans chaque pays qui déploie un réseau télégraphique~--- un fait d'observation, pas encore un mécanisme. Ce qui suit examine si cette régularité technique s'accompagne d'une régularité institutionnelle, ou si les conditions d'émergence~--- radicalement différentes d'un pays à l'autre~--- produisent des formes organisationnelles tout aussi différentes.

\dimmark{D2}{*}%
En Angleterre, la demande initiale ne vint ni de l'État ni du commerce, mais des compagnies ferroviaires, qui avaient besoin de signaler les trains sur voie unique. William Fothergill Cooke (1806--1879), ancien officier de l'armée des Indes reconverti en étudiant en médecine, assista le 22~avril 1836 à Heidelberg à une démonstration du télégraphe de Schilling~; il abandonna la médecine sur-le-champ et se consacra au télégraphe, qu'il identifiait d'emblée comme un instrument ferroviaire, non comme un moyen de communication publique \parencite{kooij_telegraph_2015, hubbard_cooke_1965}. Charles Wheatstone (1802--1875), facteur d'instruments de musique devenu expérimentateur en électricité, apportait la compétence technique qui manquait à Cooke~; ensemble, ils installèrent en 1839 une liaison d'une vingtaine de kilomètres pour la \emph{Great Western Railway} \parencite{hubbard_cooke_1965}. Le réseau se développa ensuite par l'initiative privée, sans intervention étatique\footnote{Le médecin Edward Davy (1806--1885) avait, dès 1838, démarché huit compagnies ferroviaires et créé une société, la \emph{Voltaic Telegraph Company}, capitalisée à 500\,000~livres~; il émigra en Australie avant d'aboutir, et son brevet fut racheté pour 600~livres \parencite{kooij_telegraph_2015}.}. Mais l'initiative privée ne produisit pas la concurrence~: elle produisit la concentration. \textcite{kieve_electric_1973} retrace la dynamique~: en 1852, l'\emph{Electric Telegraph Company} dominait avec environ quatre mille miles en service~; la croissance s'accompagna d'un service médiocre dans les zones peu rentables et de tarifs élevés~; la presse, particulièrement déçue, milita pour la nationalisation, qui intervint en 1869\footnote{Le \emph{Telegraph Act} de 1863 avait tenté de corriger les abus sans succès.}. L'initiative privée avait construit le réseau~; la concentration l'avait transformé en monopole~; et c'est l'État, sous la pression de la presse et non du marché, qui finit par en prendre le contrôle.

\dimmark{D2}{*}%
Aux États-Unis, le contexte était entièrement différent. Aucune compagnie privée ne portait la demande~; c'est l'État fédéral qui finança la première ligne, sur une promesse de souveraineté~--- le Congrès débattit pendant six ans de la question de savoir si la communication instantanée relevait de l'utilité générale et justifiait un financement fédéral\footnote{Le débat portait sur la souveraineté et l'utilité publique au sens politique du terme, non sur le concept de \og bien public\fg{} au sens technique de \textcite{samuelson_pure_1954} (non rival, non excluable), qui est postérieur d'un siècle.}. \textcite{thompson_wiring_1947} détaille les six années d'obstruction parlementaire~--- les représentants du Tennessee et de l'Alabama firent obstacle au vote de la subvention~--- avant que la première ligne Washington-Baltimore (64~km) n'ouvre en mai 1844. Le réseau se développa ensuite dans une concurrence anarchique~--- vingt compagnies en 1850, quatorze dans le seul Ohio en 1852, profits insignifiants, comptes mal tenus, tarifs arbitraires \parencite{bertho_histoire_1984, nonnenmacher_telegraph_2001}. La concentration capitaliste l'emporta~: le 12~juin 1866, l'\emph{American Telegraph Company} et la \emph{Western Union} fusionnèrent pour former le premier monopole privé américain \parencite{thompson_wiring_1947, nonnenmacher_telegraph_2001}. Le monopole n'avait été ni voulu ni prévu par le Congrès qui avait financé la première ligne~; il était le produit d'une concurrence que personne n'avait su réguler.

\dimmark{D2/D8}{*}%
En France, un réseau de télégraphie optique couvrait déjà le territoire, et son administration~--- militaire dans son principe, rattachée au ministère de l'Intérieur~--- n'avait aucune raison de céder la place. \textcite{ducamp_telegraphe_1869} rapporte en détail la fraude boursière qui déclencha le changement~: entre 1834 et 1836, les frères Blanc faisaient passer les cours de la rente d'État sur la ligne Paris-Bordeaux aux spéculateurs provinciaux avant l'arrivée de la malle-poste\footnote{Le procès (Tours, mars 1837) se termina par un acquittement faute de base légale. Alexandre Dumas s'en inspire dans \emph{Le Comte de Monte-Cristo} (1844), chapitres LX--LXI.}. Aucune loi n'interdisant la fraude, le ministère fit voter la loi du 2~mai 1837, dont l'article unique établissait le monopole d'État~: \og \emph{Quiconque transmettra sans autorisation des signaux d'un lieu à un autre, soit à l'aide de machines télégraphiques, soit par tout autre moyen, sera puni}\fg{}\footnote{Loi du 2~mai 1837, article unique (\emph{Bulletin des Lois}, 1837).}. Le débat avait vu s'affronter trois thèses~: liberté totale~--- les libéraux soutenaient que \og \emph{le monopole des télégraphes ne doit pas être plus concédé au gouvernement que le monopole de la presse}\fg{}\footnote{Les trois positions sont rapportées par \textcite[pp.~22--23]{bertho_histoire_1984}, qui s'appuie sur les comptes rendus de la Chambre.}~---, accès régulé au réseau d'État (les hommes d'affaires), monopole strict (le ministre). Le comte de Gasparin, ministre de l'Intérieur, l'emporta en invoquant un argument économique~--- un réseau de communication tendait par nature vers le monopole~--- et un argument policier, que le \emph{Moniteur universel} du 6~janvier 1837 rapporte en ces termes~: \og \emph{Que serait-il arrivé si les canuts insurgés de Lyon avaient pu transmettre le signal de la révolte sur tout le territoire~?}\fg{} Et Gasparin d'ajouter~: \og \emph{Les poisons et les poudres ne sont débités que par l'autorisation de l'État et certes la télégraphie entre des mains malveillantes pourrait devenir une arme dangereuse.}\fg{} Comme le documentent \textcite{ducamp_telegraphe_1869} et \textcite{belloc_telegraphie_1894}, l'administration résista ensuite au télégraphe électrique pendant douze ans~: elle commanda aux ateliers Bréguet un appareil qui reproduisait les mouvements du sémaphore optique~--- pour ne pas avoir à former les opérateurs à un nouveau code~--- et le lobby du télégraphe optique fit voter en 1842 des crédits pour un système de sémaphores nocturnes. Il fallut un changement de régime~--- l'arrivée de Louis Napoléon Bonaparte~--- pour que le réseau s'ouvre au public en 1851\footnote{Loi du 29~novembre 1850. Le ministre de l'Intérieur déclara devant les députés, en des termes qui marquaient le renversement complet de la position de 1837~: \og \emph{Il est impossible que le gouvernement français refuse de faire participer le commerce et l'industrie de notre pays aux facilités qu'offre le télégraphe électrique}\fg{} \parencite[p.~25]{bertho_histoire_1984}. Plus de six millions de francs-or furent alloués à la construction du réseau pour la seule année 1852 \parencite{ducamp_telegraphe_1869}.}. Le cas français est le seul où le monopole d'État sur la communication précéda la technologie qu'il allait gouverner~; c'est aussi le seul où une technique~--- le sémaphore optique~--- bloqua activement, par ses défenseurs institutionnels, la technique appelée à la remplacer.

\dimmark{0-GOV}{*}%
En Allemagne, le réseau fut d'emblée public. \textcite{siemens_lebenserinnerungen_1892} raconte comment, en 1842, quelques appareils de Wheatstone lui parvinrent alors qu'il était jeune officier d'artillerie~; il s'associa avec le mécanicien Halske pour présenter en 1846 un modèle dérivé du télégraphe anglais. La première ligne, Francfort-Berlin, fut construite en 1848 pour permettre à la capitale de la Prusse de suivre les travaux du Parlement fédéral \parencite[ch.~2]{headrick_invisible_1991, bertho_histoire_1984}. La ligne fut enterrée le long des voies ferrées pour la protéger des saboteurs. Le réseau se construisit ensuite sous tutelle de la poste prussienne, en même temps que l'unité politique allemande. En Allemagne, pas de résistance institutionnelle, pas de concurrence anarchique, pas de nationalisation tardive~; le réseau fut public d'emblée, et le monopole n'eut pas besoin d'émerger~--- il fut posé.

\dimmark{D2/D9/0-GOV}{+}%
Le schéma ne s'arrêtait pas à l'Europe et à l'Amérique du Nord. En Inde, l'administration coloniale britannique inaugura dès 1851 la ligne Calcutta-Diamond Harbour~; en 1865 le réseau couvrait plus de 28\,000~miles, au service du contrôle impérial \parencite[ch.~3]{headrick_invisible_1991}. Au Japon, le gouvernement Meiji fit construire la ligne Tokyo-Yokohama en 1869, dans le cadre d'un programme de modernisation étatique accéléré~: le télégraphe y fut d'emblée un monopole public au service de l'unité nationale \parencite[ch.~4]{headrick_invisible_1991}. En Chine, le réseau fut déployé à partir des années~1870, d'abord sous concession étrangère~--- la \emph{Great Northern Telegraph Company} danoise opéra les premières lignes~--- puis sous administration impériale chinoise. \textcite{hao_wiring_2022} et \textcite{gao_communication_2021} ont montré que les effets sur l'intégration des marchés du riz furent comparables à ceux observés sur le coton transatlantique. Dans ces trois cas, les conditions initiales différaient radicalement~--- colonisation, modernisation autoritaire, concession étrangère~--- mais le résultat institutionnel fut le même~: le monopole, public ou colonial.

\dimmark{D2/0-GOV}{*}%
Que sept conditions d'émergence radicalement différentes produisent le même résultat institutionnel~--- une forme ou une autre de monopole~--- n'est pas, en soi, surprenant~: la théorie économique des monopoles naturels prédit ce résultat pour toute infrastructure de réseau à coûts fixes élevés et rendements croissants. La convergence est documentée indépendamment par des sources distinctes~: \textcite{nonnenmacher_telegraph_2001} pour les États-Unis, \textcite{kieve_electric_1973} pour l'Angleterre, \textcite{bertho_histoire_1984} pour la France, \textcite{headrick_invisible_1991} pour les réseaux coloniaux. Ce que le détour par les sept cas ajoute à la théorie, c'est autre chose~: la diversité des véhicules politiques de la concentration (contrôle policier en France, souveraineté fédérale aux États-Unis, commande impériale en Inde), le fait que le contrôle politique et les rendements croissants jouent en parallèle plutôt qu'en séquence, et surtout le cas français de non-émergence~--- douze ans de blocage par l'administration du sémaphore, de 1837 à 1851~--- qui révèle que la pression vers le monopole peut retarder la transition technique autant qu'elle l'accélère.

Le mécanisme proposé pour rendre compte de cette convergence est celui que \textcite{wu2010master} a synthétisé sous le nom de \emph{master switch}~: la combinaison des rendements croissants du réseau, du désir de contrôle politique et de la logique capitalistique de consolidation exercerait une pression structurelle vers la concentration, quelle que soit la trajectoire institutionnelle initiale. La proposition de Wu est un essai d'historien-juriste, non un article académique formalisé~; elle opère par compilation interprétative de trois mécanismes économiques chacun bien établis par ailleurs dans la littérature~: les rendements croissants de réseau \parencite{katz_network_1985, economides_economics_1996}, le contrôle politique sur les infrastructures stratégiques \parencite{mazzucato_entrepreneurial_2013, headrick_invisible_1991}, et la consolidation capitalistique des infrastructures lourdes \parencite{chandler_visible_1977, perez_technological_2002}. Ce n'est donc pas l'invocation d'un théorème mais l'usage d'un cadre interprétatif qui rassemble trois résultats économiques solides en une régularité descriptive applicable aux réseaux de communication. Cette explication est compatible avec les sept trajectoires observées~; elle relève toutefois de ce que l'on pourrait appeler, en suivant la distinction de \textcite[ch.~6]{trannoy_economistes_2025}, une causalité structurelle~: un modèle théorique qui propose un mécanisme, non un résultat isolé par une variation exogène. Ce que les sept cas permettent d'établir, c'est la régularité du résultat~; ce que Wu propose, c'est un mécanisme qui en rend compte~; ce qui reste ouvert, c'est de savoir si cette pression est propre aux réseaux de communication ou commune à toutes les infrastructures lourdes.

\emergence{D2~--- Sept trajectoires institutionnelles, un seul résultat~: le monopole. La convergence est une régularité documentée, attendue par la théorie des monopoles naturels~; ce que les cas révèlent en propre, c'est la diversité des véhicules politiques, le jeu parallèle rendements croissants / contrôle étatique, et la non-émergence française comme révélateur des conditions de blocage. \questionch{2}{D2~--- Cette convergence se reproduit-elle pour le téléphone, la radio, Internet~? Est-elle un trait des réseaux de communication en tant que tels, ou des infrastructures lourdes en général~?}}

\emergence{D8~--- La loi de 1837, votée pour un réseau de sémaphores, s'applique encore au télégraphe électrique grâce à la formule \og soit par tout autre moyen\fg{}. C'est un cas où le cadre réglementaire survit à la technologie pour laquelle il a été conçu. \questionch{2}{D8~--- Ce phénomène de persistance réglementaire se retrouve-t-il aux vagues technologiques ultérieures~?}}

\emergence{D9/0-GOV~--- La géographie du déploiement télégraphique reflète les rapports de puissance~: réseau colonial en Inde, concession étrangère en Chine, domination britannique sur les câbles sous-marins. \questionch{2}{D9/0-GOV~--- La localisation des infrastructures informationnelles suit-elle la géographie du pouvoir économique et politique~?}}

La tarification du télégramme~--- au mot dans la plupart des systèmes~--- imposait une discipline de compression du langage et donnait à l'information un prix unitaire explicite. Sur les réseaux nationaux, les prix baissèrent considérablement, mais selon des mécanismes différents~: par la concurrence aux États-Unis~--- 1,55~dollar pour dix mots entre New York et Chicago en 1850, 40~cents en 1890 selon les données de \textcite{nonnenmacher_telegraph_2001}~---, par la politique tarifaire volontariste en France, où une dépêche de vingt mots de Paris à Marseille coûtait 18,44~francs en 1851 et 1~franc en 1878 \parencite{ducamp_telegraphe_1869}\footnote{Le tarif de 1851 représentait, selon la formule de \textcite{bertho_histoire_1984}, pp.~33--35, \og le prix d'un vêtement du dimanche qu'une ouvrière portera toute l'année\fg{}.}. Le câble transatlantique, en revanche, restait réservé au trafic d'affaires~: \dimmark{D5}{*}%
10~dollars par mot en 1866\footnote{Le bureau de Londres de l'\emph{Associated Press} envoyait 35\,000~mots par jour en 1869, pour des dépenses atteignant 2\,000~dollars par jour \parencite{frederix_siecle_1959, bertho_histoire_1984}.}.

\paragraph{Le câble transatlantique et la formation des marchés}\mbox{}\\
La coordination marchande, elle, dispose d'un résultat d'une force probante rare dans cette section, fondé sur une stratégie d'identification économétrique. Avant le câble transatlantique de 1866, le prix du coton à Liverpool et celui de New York pouvaient diverger pendant les dix jours de traversée~; la bourse du coton de New York fut créée en 1870, quatre ans après la pose du câble \parencite{standage1998victorian, gordon_thread_2002}. \textcite{steinwender_real_2018}, à partir de données quotidiennes du \emph{New York Times} et du \emph{Liverpool Mercury} et d'une stratégie d'identification fondée sur la rupture temporelle du câble, a montré que la connexion transatlantique réduisit l'écart de prix moyen de 35\,\%, l'équivalent de la suppression d'un tarif \emph{ad valorem} de 7\,\%, et que les gains d'efficience représentaient environ 8\,\% de la valeur annuelle des exportations. L'effet n'était pas symétrique~: les prix à New York réagissaient fortement aux nouvelles de Liverpool, tandis que Liverpool réagissait peu aux nouvelles de New York~--- parce que le marché britannique était le marché dominant du coton mondial, et que New York était preneur de prix, non faiseur. Surtout, l'effet n'était pas seulement redistributif~: les exportations elles-mêmes répondaient aux nouvelles de prix, ce qui signifie que le câble ne se contentait pas de réallouer les profits entre acteurs mieux et moins bien informés~--- \dimmark{D6}{*}%
il créait du commerce qui n'aurait pas eu lieu sans lui. Un effet comparable avait déjà été observé sur les marchés intérieurs américains~: en 1846, les prix du blé et du maïs à Buffalo avaient quatre jours de retard sur ceux de New York~; en 1848, après la connexion télégraphique entre les deux villes, les prix étaient synchronisés \parencite{nonnenmacher_telegraph_2001}.

\textcite{juhasz_spinning_2018} ont montré, sur les textiles de coton britanniques exportés vers 75~pays entre 1845 et 1880, que l'effet du télégraphe dépendait de la \emph{codifiabilité} du produit~: le fil (\emph{yarn}), dont la variété se résumait à un chiffre standardisé, vit ses importations augmenter de 10,6\,\% après la connexion télégraphique~; le tissu uni (\emph{plain cloth}), dont les caractéristiques pouvaient être décrites en quelques mots, de 5,0\,\%~; le tissu imprimé (\emph{finished cloth}), dont les motifs ne pouvaient être spécifiés que par l'examen d'un échantillon physique, ne bénéficia d'aucun effet statistiquement significatif\footnote{Les élasticités estimées par variables instrumentales (la rugosité du fond marin comme prédicteur de la date de connexion) sont de $-0{,}27$ pour le fil, $-0{,}13$ pour le tissu uni, et $-0{,}03$ (non significatif) pour le tissu fini. \textcite{juhasz_spinning_2018} montrent également que le télégraphe facilita la diffusion technologique internationale, les pays connectés important davantage de machines et d'intrants intermédiaires~--- l'ICT ne facilitait pas seulement le commerce, elle accélérait le transfert de savoir technique.}. \dimmark{D1}{?}%
Le télégraphe accélérait les mots, pas les choses~: pour les produits codifiables, il ne réduisait pas un coût de transaction préexistant mais \emph{rendait possible la spécification à distance}, et donc la commande, qui n'aurait pas eu lieu sans lui~; pour les produits non codifiables, l'échantillon continuait de voyager par bateau, et rien ne changeait. On notera l'asymétrie~: la ligne télégraphique est un bien rival~--- un message occupe le circuit, et le suivant doit attendre~---, mais l'information qu'elle transporte ne l'est pas~--- une fois le cours du coton connu à Liverpool, rien n'empêche sa diffusion immédiate à tous les acteurs du marché.

\emergence{D6/D9/D10~--- Le résultat de \textcite{juhasz_spinning_2018} soulève trois questions structurelles. Premièrement, l'effet de l'ICT sur le commerce n'est pas uniforme~: il dépend de la correspondance entre ce que la technologie sait transmettre et ce que le produit exige comme spécification. Si ce résultat se transpose, chaque saut technologique~--- du télégraphe (mots) au téléphone (voix) à Internet (images, vidéo)~--- redéfinirait la frontière entre les biens spécifiables à distance et ceux qui exigent la co-présence physique. \questionch{2}{Cette frontière est-elle la même que celle de Rauch (1999) entre biens \og organisés\fg{} et biens \og différenciés\fg{}~? L'ICT la déplace-t-elle continûment~?} Deuxièmement, le télégraphe a davantage augmenté le commerce d'intrants intermédiaires (le fil, $+10{,}6\,\%$) que de produits finis (le tissu imprimé, effet nul)~: le télégraphe ne se contente pas d'augmenter le volume du commerce, il en modifie la \emph{composition}, en favorisant la fragmentation de la chaîne de valeur. Un pays connecté importe du fil au lieu du tissu fini~--- et se met à tisser localement, important aussi des machines. Le télégraphe facilite le commerce, mais il réorganise aussi la division internationale du travail et accélère le transfert technologique. \questionch{2}{Ce mécanisme de fragmentation~--- que la littérature contemporaine documente pour Internet et l'offshoring \parencite{grossman_trading_2008}~--- opérait-il déjà au \textsc{xix}\ieme{} siècle avec le télégraphe~? Si oui, la codifiabilité est moins une dimension parmi d'autres qu'un filtre qui module l'effet de toutes les dimensions.} Troisièmement, la codifiabilité n'est pas une propriété fixe des produits~--- elle est relative à la technologie disponible. Le tissu imprimé qui résistait au télégraphe (il fallait voir l'échantillon) deviendrait potentiellement codifiable dès qu'une technologie saurait transmettre des images. Si cette logique se poursuit, chaque saut technologique ne se contenterait pas d'accélérer la transmission de ce qui est déjà codifiable~: il \emph{rendrait codifiable ce qui ne l'était pas}, et déplacerait la frontière entre les transactions réalisables à distance et celles qui exigent la co-présence. \questionch{2}{Si l'IA générative rend codifiable presque tout (image, voix, intention, spécification), est-elle une ICT de plus dans la série, ou une ICT qui élimine le filtre même qui limitait toutes les précédentes~?} \questionch{3}{Si la codifiabilité croissante fragmente les chaînes de valeur de services comme le télégraphe a fragmenté celle des textiles, chaque fragment échangé à distance consomme de la bande passante, du calcul et du stockage~--- c'est-à-dire de l'énergie. La fragmentation informationnelle est-elle un vecteur de rebond énergétique~?}}

Le résultat de \textcite{steinwender_real_2018} n'est pas isolé~: il a été répliqué sur d'autres produits, dans d'autres géographies, par des méthodes différentes. \textcite{ejrnaes_gains_2010} montrent sur le blé Chicago-Liverpool que la demi-vie des chocs à la loi du prix unique passe de 3,5 à 2,6~semaines après le câble, et que les gains d'efficacité du marché sont d'un ordre de grandeur comparable aux gains liés à la baisse des coûts de transport~--- l'information n'est pas un facteur secondaire. \textcite{hao_wiring_2022} répliquent le résultat dans une économie radicalement différente~--- le riz chinois, neuf provinces, 1870-1911 ~--- et observent en outre que la volatilité des prix \emph{augmente} dans les marchés auparavant isolés, exactement comme les flux transatlantiques de Steinwender, parce que la connexion informationnelle expose les marchés locaux à des chocs distants qu'ils ne percevaient pas. \textcite{andrabi_information_2021}, sur le grain indien (192~districts, 1862-1920), isolent l'effet informationnel de l'effet transport en montrant que le télégraphe réduit la dispersion des prix \emph{avant} l'arrivée du chemin de fer ~--- mais seulement là où un moyen de transport physique existe déjà (la Gange, la Grand Trunk Road). L'information sans transport ne produit pas d'arbitrage\footnote{\textcite{gao_communication_2021}, dans un résultat complémentaire sur le réseau chinois~--- où les chemins de fer sont quasi absents ---, montrent que le télégraphe réduit l'incidence des prix extrêmes et atténue la réponse des prix aux chocs météo locaux, mais augmente la réactivité aux chocs dans les régions connectées à distance. C'est un résultat de \emph{risk-sharing} spatial~: la volatilité converge~--- les marchés stables devenant plus volatils, les marchés instables se stabilisant. L'effet dépend de la densité de voies navigables, ce qui confirme la conditionnalité au transport.}. \dimmark{D6}{+}%
L'ensemble dessine une régularité convergente~: le télégraphe modifie la coordination marchande partout où il est déployé, mais cet effet est conditionnel à l'existence d'un moyen d'arbitrage physique, et il a un coût~--- l'exposition des marchés locaux à la contagion des chocs distants.

\emergence{D6~--- Le télégraphe ne réduit pas un coût de transaction préexistant~: il rend possible la transaction elle-même. Le câble transatlantique crée du commerce qui n'aurait pas eu lieu sans lui \parencite{steinwender_real_2018}, et cet effet est répliqué sur le blé, le riz chinois et le grain indien. Mais il est conditionnel~: sans moyen de transport physique, l'information ne produit pas d'arbitrage~; et il a un coût~: l'exposition des marchés locaux à la contagion des chocs distants. \questionch{2}{D6~--- Ce rôle constitutif de l'information synchrone se retrouve-t-il dans les marchés coordonnés par le téléphone, puis par Internet~?}}

L'augmentation du trafic se heurta aux capacités des lignes. Selon \textcite{beauchamp2001telegraphy}, les appareils Morse transmettaient 25~mots à la minute~; les appareils Hughes, déployés en France à partir de 1862, portèrent ce débit à 45~mots~; le système Baudot (1874), fondé sur le principe du temps partagé, atteignit 60~mots. Mais ces gains de débit étaient conditionnels à la qualité matérielle du réseau. \textcite{brooks_report_1875}, ingénieur télégraphique américain envoyé à l'Exposition de Vienne en 1873, observait que le Hughes français tournait sous la pluie à pleine vitesse sur Marseille-Paris ou Lyon-Paris~; aux États-Unis, sur des lignes mal isolées et hâtivement posées, la même technologie était inutilisable, et les opérateurs revinrent à la lecture par sondeur du Morse, plus rustique mais robuste à l'instabilité électrique. \dimmark{0-TER/D9}{*}%
La matérialité du substrat ne conditionnait pas seulement l'extraction des inputs~; elle filtrait quels artefacts pouvaient effectivement fonctionner sur le réseau, et donc l'horizon des trajectoires techniques nationales. Parallèlement, le multiplexage progressait~: duplex de Stearn (1870), quadruplex d'Edison (1874)\footnote{Thomas Alva Edison (1847--1931), ancien opérateur de télégraphe itinérant sur les lignes du Midwest, avait passé ses nuits à expérimenter avec les circuits. Le quadruplex~--- quatre messages simultanés sur un seul fil~--- fut sa première grande invention et lui rapporta 100\,000~dollars de Jay Gould et de Western Union \parencite{kooij_telegraph_2015}. C'est cette fortune et cette maîtrise du circuit électrique, acquises dans le domaine de la communication, qu'il transposa dans le domaine de la puissance~: le quadruplex est le pont entre le monde du télégraphe et celui de l'électrification, via un homme qui était dans les deux. Les capacités transatlantiques passèrent de 15--17~mots par minute vers 1870 à 90~mots par minute après l'introduction de câbles de gros diamètre en 1894.}. Mais c'est en cherchant à réaliser le multiplexage par différenciation des fréquences\footnote{Le multiplexage par différenciation des fréquences consiste à transmettre simultanément plusieurs signaux sur un même fil en attribuant à chacun une fréquence porteuse distincte. Le récepteur sépare les signaux en filtrant les fréquences. C'est le principe qui sera repris, un siècle plus tard, par la radio FM et les réseaux cellulaires.}~--- le \og télégraphe harmonique\fg{}~--- que Graham Bell et Elisha Gray découvrirent le téléphone. Alexander Graham Bell (1847--1922), enseignant pour sourds à Boston, fils d'Alexander Melville Bell (inventeur du \emph{Visible Speech}), expérimentait le soir dans son atelier avec des électro-aimants, des bobines et des diapasons. Il était financé par Gardiner Hubbard, père d'une de ses élèves sourdes~--- Mabel, qu'il épousera~--- et par George Sanders, père d'un autre. Hubbard voulait briser le monopole de Western Union~; il finançait Bell pour trouver une alternative au quadruplex d'Edison. C'est en travaillant sur ce problème commercial~--- le multiplexage~--- que Bell découvrit la transmission de la voix \parencite{kooij_telephone_2016}. Elisha Gray (1835--1901), ingénieur à Western Electric, déposa son \emph{caveat} pour le téléphone le même jour que Bell~--- le 14~février 1876. Gray est l'homme qui ne vit pas la valeur du téléphone parce que sa métrique était celle du télégraphe~: messages par fil, coût par mot. Le téléphone était un sous-produit de la compétition télégraphique, engendré par la recherche d'économies de cuivre. \dimmark{D2/D10}{+}%
La trajectoire du multiplexage donna ainsi naissance à deux inventions qui la dépassaient~: le téléphone de Bell et, par Edison, l'électrification.

\emergence{Le multiplexage est D2 par excellence~: le coût de la ligne est fixe, toute technique qui augmente le nombre de messages par fil augmente le rendement. \questionch{2}{D10~--- la course au débit est-elle une constante structurelle, de Baudot à la 5G~?}}

La guerre de Crimée (1853--1856) illustre un mécanisme d'émergence distinct de tous les précédents~: la demande militaire, financée sans contrainte budgétaire par l'État en guerre, créa en quelques mois une ligne Paris-Vienne-Bucarest-Crimée que le commerce n'aurait pas justifiée. Le télégraphe transforma les conditions de commandement~--- au point que le général Pélissier demanda à être relevé d'un commandement qu'il jugeait \og \emph{\dimmark{D7/0-GOV}{*}%
impossible à exercer à l'extrémité parfois paralysante d'un fil électrique}\fg{}\footnote{Échange entre Napoléon~III et Pélissier rapporté dans \textcite{frederix_siecle_1959}, \textcite{ducamp_telegraphe_1869} et \textcite[p.~44]{bertho_histoire_1984}. L'expression de Pélissier (1855) est peut-être la première formulation du problème de la micro-gestion à distance par les télécommunications. L'agence Havas avait célébré à tort la prise de Sébastopol en septembre~1853, ce qui faillit ruiner son crédit \parencite{frederix_siecle_1959}.}.

Ce que ces usages militaires, commerciaux et financiers ne comptabilisaient pas, c'est ce que le signal présupposait en amont de lui-même. Le signal télégraphique consommait une quantité d'énergie négligeable~--- les piles de Daniell fournissaient quelques milliampères dont le coût n'apparaît nulle part dans les comptes d'exploitation\footnote{Argument \emph{ex silentio}. Western Union dégageait 30 à 40\,\% de marge nette \parencite{nonnenmacher_telegraph_2001}~; les comptes d'exploitation publiés ne détaillent pas les postes de consommation énergétique. \textcite{kieve_electric_1973} fournit des bilans comptables pour le réseau britannique, mais sans décomposition des coûts par nature de ressource.}. \dimmark{0-TER}{+}%
Mais le système reposait sur une chaîne d'extraction dont l'ampleur n'a été documentée que rétrospectivement.

\dimmark{D9/0-TER}{+}%
Aux États-Unis, le réseau terrestre exigeait des centaines de milliers de poteaux en bois dont l'approvisionnement impliquait, comme l'a montré \textcite{fitzmaurice_material_2023}, une déforestation géographiquement délocalisée. Western Union qualifiait ses poteaux de \emph{perishable property}~; les poteaux non traités duraient environ dix ans \parencite{gordon_capital_nineteenth}. Le réseau terrestre laissait entrevoir ce que les câbles sous-marins allaient confirmer à plus grande échelle~: l'infrastructure de transmission avait une empreinte matérielle dont l'ampleur n'apparaissait dans aucun bilan.

\dimmark{D9/0-GOV}{*}%
Sur les câbles sous-marins, en effet, la dépendance matérielle était plus spectaculaire encore. Entre 1850 et 1880, comme le documentent \textcite{headrick_invisible_1991}, \textcite{hunt_imperial_2021} et \textcite{bertho_histoire_1984}, la télégraphie sous-marine internationale était entièrement dominée par l'Angleterre~: les financiers de Londres contrôlaient l'approvisionnement en cuivre et en \emph{gutta-percha}, fixaient les prix et les quantités\footnote{La concentration s'étendait à la fabrication elle-même~: \textcite{muller_telegraph_2015} documentent que Telcon (\emph{Telegraph Construction and Maintenance Company}) manufacturait sur les rives de la Tamise la quasi-totalité des câbles sous-marins mondiaux jusque dans les années~1930.}. Le câble de 1865 pesait 980~kg/km, dont 73~kg/km de cuivre et 98~kg/km de gutta-percha~--- un latex naturel extrait d'arbres tropicaux du genre \emph{Palaquium}, endémiques de Malaisie et Bornéo, dont un arbre mature de trente ans produisait environ 600~grammes utilisables \parencite{greene_dark_2005, sanzo_around_2023}. \textcite{tully_victorian_2009} a estimé, par une reconstruction rétrospective fondée sur les tonnages d'exportation, que les seuls câbles atlantiques de 1865--66 avaient nécessité l'abattage d'au moins 1,2~million d'arbres\footnote{Estimation reconstruite, non décompte exhaustif.}. Telcon estime avoir consommé 47~millions d'arbres sur un siècle \parencite{greene_dark_2005}. Les prix grimpèrent de 8 à 60~dollars espagnols par picul\footnote{Le picul (équivalent à environ 60~kg) était l'unité de mesure standard du commerce de la gutta-percha à Singapour. Le dollar espagnol, monnaie de référence du commerce asiatique au \textsc{xix}\ieme{} siècle, valait environ cinq francs-or ou un dollar américain. La hausse de 8 à 60~dollars par picul en moins de dix ans représente donc une multiplication par sept et demi du prix de la matière première.} entre 1844 et 1853~; les arbres furent épuisés à Singapour dès 1847 \parencite{sanzo_around_2023, buckley_anecdotal_1902}.

\dimmark{D8/D9}{*}%
Le cas du câble transatlantique révèle, au-delà de la matérialité, la structure du financement et la promesse qui le porte. Cyrus Field, homme d'affaires new-yorkais sans compétence technique, finança l'entreprise sur une promesse de rente informationnelle~: celui qui contrôlerait le câble contrôlerait l'arbitrage entre les marchés du coton, du blé et des titres. La promesse mobilisa des capitaux britanniques considérables~--- mais l'accès aux matières premières (cuivre, gutta-percha) constituait une contrainte que personne ne quantifiait, et le financement lui-même était une aventure. Le premier câble (1857--1858), qui ne fonctionna que trois semaines, avait coûté environ 300\,000~livres sterling à l'\emph{Atlantic Telegraph Company}~; la compagnie fit faillite et ne fut reconstituée qu'en 1864, avec les capitaux de Pender et Brassey, pour un câble réussi en 1866 dont le coût atteignit 600\,000 à 700\,000~livres \parencite{gordon_thread_2002}\footnote{Le pattern investissement massif en infrastructure ICT $\rightarrow$ faillite $\rightarrow$ recapitalisation $\rightarrow$ succès se reproduit avec la \emph{railway mania} des années~1840, la \emph{corporate mania} électrique des années~1880--1890 et la bulle dot-com de~2000.}. Les faillites des compagnies françaises de câbles sous-marins~--- la \emph{Compagnie Erlanger} rachetée en 1873, la \emph{Compagnie française du câble Paris-New York} en faillite en 1895~--- montrèrent les limites du développement par concessions dans un secteur à capitaux massifs et risque immense \parencite{bata_cables_1982, headrick_invisible_1991}. Aux faillites par concession s'ajoute, à partir des années~1890, un autre type d'échec~: l'infrastructure posée \emph{avant} la demande. \textcite{muller_telegraph_2015} documentent comment les câbles transpacifiques de 1902 (canadien) et 1904 (américain), reliant l'Amérique du Nord à l'Asie via Hawaï et Midway, échouèrent commercialement. Le trafic ne s'approcha jamais des coûts d'infrastructure. L'Atlantique de 1866 avait suivi une route commerciale dense préexistante~; les câbles pacifiques, eux, furent posés pour ouvrir un marché que l'on espérait voir naître~: la rivalité sino-américaine, les ambitions canadiennes vers l'Asie, la doctrine de la porte ouverte fournissaient les motifs~; la demande commerciale, elle, n'existait pas. \dimmark{D6/D8}{*}%
C'est la position symétrique du blocage français de 1837-1851~: là, la demande existait mais l'infrastructure ne se construisait pas~; ici, l'infrastructure se construit mais la demande ne se manifeste pas. Le câble pacifique fonctionne, le réseau est complet, et les comptes ne rentrent pas. Le réseau mondial reposait sur du cuivre des Cornouailles et du Michigan, du bois du Wisconsin, du latex de Bornéo, du charbon des navires câbliers. \textcite{belloc_telegraphie_1894}, qui écrivait au c\oe{}ur de cette période, constatait dans sa préface que la télégraphie électrique avait été \og surtout remarquable par ses conséquences économiques\fg{}~--- sans éprouver le besoin de mentionner ses conséquences matérielles. La coïncidence est au moins suggestive~: en 1865, la même année que le premier câble transatlantique, \textcite{jevons_coal_1865} publiait \emph{The Coal Question}. Les préoccupations énergétiques de l'Angleterre victorienne portaient sur l'épuisement du charbon et l'avenir de la puissance industrielle~--- pas sur le cuivre, le latex ou le bois que le réseau de communication consommait. On ne peut pas en conclure que l'occultation était structurelle, mais on peut constater que les catégories de l'époque séparaient l'énergie~--- question de puissance~--- et l'information~--- question de commerce et de contrôle. Cette séparation, qu'elle fût délibérée ou simplement un effet des ordres de grandeur, n'a pas été corrigée depuis~: ni les historiens de l'énergie \parencite{smil2017energy, wrigley2010energy}, ni les historiens de l'ICT \parencite{headrick_invisible_1991, standage1998victorian}, ni les \emph{environmental historians} n'ont fait de la relation entre infrastructure informationnelle et substrat matériel un objet d'étude dans la longue durée\footnote{Les seules exceptions sont \textcite{tully_victorian_2009} sur la gutta-percha, \textcite{fitzmaurice_material_2023} sur la matérialité du réseau américain, et l'article programmatique de \textcite{ensmenger_environmental_2018} sur l'histoire environnementale du \emph{computing}. \textcite{magalhaes_accumuler_2024} a produit pour les infrastructures de transport françaises (routes, autoroutes, béton) le type d'histoire environnementale qui manque pour les réseaux de communication. Il n'existe à ce jour aucune analyse du cycle de vie du réseau télégraphique du \textsc{xix}\ieme{} siècle au sens ISO~14040/44.}.

\questionch{3}{Matérialité occultée~: le pattern gutta-percha/poteaux/cuivre préfigure-t-il les terres rares, le lithium et l'eau de refroidissement des datacenters~? La domination britannique sur les matériaux du câble préfigure-t-elle la domination américaine sur les semi-conducteurs ou chinoise sur les terres rares~? \textcite{pirson_environmental_2022} montrent que l'empreinte environnementale par cm\textsuperscript{2} de fabrication de circuits intégrés, qui baissait depuis les années~1990, remonte à partir du n\oe{}ud de 7~nm~: la miniaturisation exige un bouquet élémentaire toujours plus large et plus pur, dont les impacts croissent avec l'ambition fonctionnelle~--- le même pattern, à un siècle et demi de distance.}

\emergence{D1/D6/0-TER~--- Steinwender montre que le stockage physique de coton atténue les frictions informationnelles mais ne les élimine pas~: les négociants utilisent leurs stocks comme tampon face à l'incertitude sur les prix distants. L'analogie avec le stockage par batteries dans la transition énergétique est à analyser~: la batterie atténue l'intermittence des renouvelables mais ne l'élimine pas. Dans les deux cas, l'information en temps réel (le câble) se \emph{substitue} partiellement au stock physique~; si l'analogie tient, le smart grid pourrait jouer un rôle comparable vis-à-vis des batteries~--- le couple information + stockage serait plus puissant que chacun des deux seuls. \questionch{3}{La complémentarité entre fonctions ICT (transmit = câble, store = stock) et substrat matériel est-elle une constante structurelle du couplage ICT-énergie~?}}

Avant de poursuivre, un point de pause~--- non sur le contenu mais sur le statut de ce qui vient d'être avancé. Le récit qui précède a procédé par trois voies que les historiens et les économistes ne confondent pas toujours, et qu'il importe de distinguer\footnote{La distinction entre ces trois registres s'inspire des catégories proposées par \textcite{monnet_perspectives_2025} dans sa réflexion sur les discordes méthodologiques entre économistes et historiens~: causalité processuelle (le récit de l'enchaînement des actions), causalité structurelle (un modèle théorique qui propose un mécanisme), causalité interventionniste (une quasi-expérience qui estime un effet). Voir \textcite[ch.~6--8]{trannoy_economistes_2025} pour un traitement approfondi de ces questions.}.

La reconstruction narrative des trajectoires nationales~--- qui a déployer, pour quelle demande, dans quel cadre institutionnel, et avec quel résultat~--- relève de ce que l'on peut appeler une évidence processuelle~: un enchaînement documenté d'actions, de décisions et de conséquences, attesté par des sources indépendantes. La substitution du capital au travail dans sept contextes institutionnels \parencite{gordon_capital_nineteenth, holzmann_early_1995, gabler_american_1988}, la convergence vers le monopole \parencite{wu2010master, nonnenmacher_telegraph_2001, kieve_electric_1973, bertho_histoire_1984}, le pattern investissement-faillite-recapitalisation du câble \parencite{gordon_thread_2002, headrick_invisible_1991}, la matérialité extractive du réseau \parencite{fitzmaurice_material_2023, tully_victorian_2009}~: ces régularités sont documentées indépendamment par des sources distinctes. Elles ne dépendent pas d'un cadre théorique pour être observées.

Les mécanismes proposés pour en rendre compte relèvent d'un autre registre. Le \emph{master switch} de Wu (rendements croissants + contrôle politique + consolidation), la codifiabilité de Juhász comme filtre des effets du commerce, le seuil de McCallum au-delà duquel la coordination hiérarchique devient nécessaire~: chacun de ces mécanismes est compatible avec le matériau observé, mais aucun n'est isolé par une variation exogène. Ce sont des propositions explicatives, non des résultats testés.

Le résultat de \textcite{steinwender_real_2018} relève d'un troisième registre encore~: une stratégie d'identification économétrique fondée sur la rupture temporelle du câble transatlantique, qui permet d'estimer l'effet de l'information synchrone sur l'intégration des marchés. Ce résultat a été répliqué dans quatre autres géographies par des méthodes différentes \parencite{ejrnaes_gains_2010, hao_wiring_2022, andrabi_information_2021, gao_communication_2021}, ce qui renforce sa portée~; il reste toutefois conditionnel aux hypothèses de l'instrument et au contexte spatio-temporel dans lequel il a été estimé, et ne se transforme pas en loi universelle.

Parmi les pistes qui restent au stade de l'hypothèse~: la tension entre efficacité informationnelle et empreinte matérielle, la substitution partielle entre information et stockage physique, et la question de savoir si l'information synchrone modifie les préférences ou seulement les contraintes.

Mais le télégraphe ne fonctionnait jamais seul. Il s'insérait dans des systèmes techniques plus vastes, et c'est dans cette insertion que l'on peut observer comment une infrastructure de transmission interagit avec un substrat physique et énergétique concret. Le cas qui suit~--- le télégraphe dans le chemin de fer~--- n'est pas un septième cas d'émergence nationale~; c'est un changement de registre. Les sept trajectoires qui précèdent examinaient le télégraphe comme réseau autonome~: qui le déploie, sous quelle forme institutionnelle, avec quel effet sur les marchés. Le cas ferroviaire examine le télégraphe comme composant d'un système plus large~: comment une infrastructure de transmission s'insère dans un système de transport, quelles formes organisationnelles cette insertion engendre, et quel rapport le circuit informationnel entretient avec le substrat énergétique du système qu'il coordonne.


\subsection*{Le télégraphe dans le chemin de fer}\label{subsec:railroad}

\paragraph{La collision du Western Railroad (1841)}\mbox{}\\
C'est à partir du réseau ferroviaire américain et britannique, entre 1840 et 1907, que \textcite{chandler_visible_1977} a construit sa théorie de la firme multidivisionnelle et \textcite{beniger_control_1986} sa théorie de la société de l'information. Le phénomène n'était pas limité aux pays anglo-saxons~--- les réseaux prussiens, le PLM français\footnote{La Compagnie des chemins de fer de Paris à Lyon et à la Méditerranée, fondée en 1857, exploitait le plus grand réseau français avec plus de 10\,000~km de lignes et utilisait le télégraphe pour la coordination du trafic dès les années~1860.}, le réseau colonial indien connaissaient des problèmes de coordination comparables~--- mais c'est le cas américain que Chandler et Beniger ont analysé avec une profondeur qui n'a pas d'équivalent ailleurs, et c'est sur cette analyse que nous nous appuyons. Le revisiter, avec les corrections que lui a apportées \textcite{schwantes_train_2019}, permet de voir non seulement ce que le télégraphe faisait dans ce système, mais aussi ce qu'il ne faisait pas~--- et surtout ce que le système tout entier impliquait comme rapport à l'énergie sans que personne, à l'époque ni depuis, ne le comptabilise.

Le problème se formulait simplement. L'application de la vapeur au transport à partir des années~1830 avait accéléré la vitesse et le volume des flux matériels dans des proportions que les dispositifs de coordination existants ne pouvaient pas absorber. \textcite{beniger_control_1986} donne à ce phénomène le nom de \emph{crise de contrôle}~: pour la première fois dans l'histoire, \dimmark{D6/D7}{*}%
des biens se déplaçaient plus vite que l'information les concernant. Le cas fondateur est celui du \emph{Western Railroad}, Massachusetts. Le 5~octobre~1841, sur 156~miles de voie unique entre Worcester et Albany, deux trains de passagers entrèrent en collision frontale~: deux morts, dix-sept blessés. L'enquête, reconstituée par \textcite[p.~183]{salsbury_state_1967}, identifia la cause~: six trains circulaient quotidiennement, nécessitant neuf croisements sur voie unique, sans signaux standardisés ni communication en temps réel. La programmation du trafic reposait sur des horaires imprimés dont l'exécution supposait une ponctualité parfaite~--- condition que la vapeur, par sa puissance même, rendait impossible à garantir\footnote{La ponctualité parfaite est une hypothèse de l'horaire~; la vapeur est une réalité physique dont la variabilité défie l'horaire. La contradiction est structurelle, pas accidentelle.}.

La crise de contrôle n'était pas un accident local. Elle se reproduisait partout où le réseau s'étendait au-delà de la portée de la supervision directe. Un surintendant de ligne, écrivait Daniel McCallum\footnote{Daniel Craig McCallum (1815--1878), ingénieur et administrateur d'origine écossaise. L'\emph{Erie Railroad}, achevé en 1851, reliait l'Hudson au lac Érié sur 483~miles à travers le sud de l'État de New York~--- à l'époque le plus long chemin de fer du monde.} en 1856 dans son rapport au conseil d'administration de l'\emph{Erie Railroad}, pouvait donner son attention personnelle à une ligne de cinquante miles~; mais à cinq cents miles, les mêmes méthodes devenaient \og \emph{entirely inadequate}\fg{} \parencite[p.~21--22]{chandler_strategy_1962}. Ce n'était pas une question de degré~: c'était un seuil qualitatif, au-delà duquel le fonctionnement du système exigeait non pas plus de la même chose, mais autre chose. Aucune organisation antérieure, notait \textcite[p.~94]{chandler_visible_1977}, n'avait nécessité une coordination aussi constante sur des centaines de miles~--- ni les armées, ni les manufactures, ni les administrations d'État.

\paragraph{McCallum, les six formes informationnelles et l'Erie Railroad}

La réponse à la crise ne fut pas le télégraphe. Elle fut un ensemble ordonné de formes informationnelles, déployées séquentiellement, dont le télégraphe n'était que la quatrième. En reprenant la chronologie établie par \textcite[ch.~5--6]{beniger_control_1986} et \textcite[p.~30, p.~148]{chandler_henry_1956}, on distingue six formes successives. Les \emph{rule books} et les horaires, d'abord, codifiaient les règles d'exploitation dans un document distribué à chaque agent~: information figée, exécutée mécaniquement, qui fonctionnait tant que les conditions réelles ne déviaient pas du plan. La hiérarchie managériale de McCallum, ensuite, instaurée sur l'\emph{Erie Railroad} à partir de 1854, créait le premier système hiérarchique formel de collecte, de traitement et de communication de l'information~: un organigramme en arbre, des uniformes codés par subdivision et par grade, une architecture de traitement~--- pas un simple signal\footnote{\textcite{yates_control_1989} montre que ce \og software organisationnel\fg{} comprenait aussi des formulaires standardisés, un classement vertical des documents, et des technologies de duplication qui rendaient l'information reproductible à l'intérieur de la firme.}. Le reporting systématique, troisièmement, rendait le réseau \emph{observable} depuis le centre~: rapports horaires, quotidiens, mensuels, compilés par huit clercs à temps plein \parencite[p.~148]{chandler_henry_1956}. Le télégraphe, quatrièmement, ajoutait la communication en temps réel~: le premier usage opérationnel documenté, sur l'\emph{Erie}, date de 1851, lorsque le surintendant Minot\footnote{Charles Minot (1810--1866), surintendant de la division de l'\emph{Erie Railroad}. L'épisode de 1851 est le premier usage documenté du télégraphe pour le contrôle opérationnel du trafic ferroviaire.} retarda un train à quatorze miles de distance par un message télégraphique \parencite{schwantes_train_2019}. La comptabilité analytique, cinquièmement, transformait les données brutes en indicateurs décisionnels~: Thomson\footnote{John Edgar Thomson (1808--1874), président du \emph{Pennsylvania Railroad} de 1852 à 1874. Le \emph{Pennsylvania} était le principal concurrent de l'\emph{Erie} et devint, sous sa direction, la plus grande entreprise privée du monde.}, sur le \emph{Pennsylvania Railroad}, introduisit le coût par tonne-mile~; Fink\footnote{Albert Fink (1827--1897), ingénieur d'origine allemande, vice-président du \emph{Louisville and Nashville Railroad}. Sa comptabilité analytique décomposait les coûts d'exploitation par segment de ligne et par type de service.} développa une comptabilité analytique complète \parencite[p.~231--232]{beniger_control_1986}\footnote{Mais \textcite{thompson_myth_1991} montre que les hypothèses de coûts fixes sur lesquelles reposait cette comptabilité étaient erronées~--- preuve que l'on peut innover dans la forme organisationnelle tout en se trompant sur le contenu analytique.}. Enfin, la standardisation infrastructurelle inscrivait l'information dans la matière elle-même~: écartement national en 1886, wagons standardisés en 1867, signaux et \emph{interlocking}\footnote{L'\emph{interlocking} est un dispositif mécanique qui rend physiquement impossible l'activation simultanée de signaux contradictoires~--- par exemple, un signal de voie libre et un aiguillage dirigeant un train vers une voie occupée. L'information de sécurité est inscrite dans le mécanisme lui-même.} en 1875, fuseaux horaires en 1883 \parencite[p.~235--236]{beniger_control_1986}.

\dimmark{D3/D10}{+}%
Ces six formes ne sont pas du même ordre. Les unes~--- le \emph{rule book}, la hiérarchie managériale, le reporting, la comptabilité analytique~--- sont des technologies organisationnelles dont le coût de diffusion est faible~: du \emph{software} social, reproductible au prix de l'impression et de la formation. Les autres~--- le télégraphe et la standardisation infrastructurelle~--- sont du \emph{hardware} matériel dont les coûts fixes sont élevés. \textcite{kooij_telegraph_2015}, dans son analyse micro-économique de l'innovation dans les réseaux de communication, formalise cette distinction sous le nom de \emph{communication engines}~: des artefacts techniques dont l'adoption exige un investissement en capital que toutes les compagnies ne pouvaient pas consentir. L'installation de la ligne télégraphique sur le \emph{Great Western Railway} coûtait 3\,270~livres sterling pour treize miles~; un \emph{rule book} ne coûtait que le prix de l'impression. La capitalisation de la compagnie devenait ainsi une variable médiatrice de l'adoption~: les moins capitalisées pouvaient s'organiser mieux, pas s'équiper. \textcite{schwantes_train_2019} l'a confirmé en corrigeant le récit de Beniger~: même en 1861, seules les grandes lignes utilisaient le télégraphe~; l'adoption massive vint après la Guerre civile, portée par des chocs exogènes~; et le \emph{Hours of Service Act} de 1907 accéléra \dimmark{D7/D4}{*}%
le passage au téléphone.

\paragraph{Le block system britannique}

Le cas britannique renversait la perspective. En Amérique, le problème était la distance~--- cinq cents miles de voie unique en territoire continental~--- et le télégraphe avait été un \og mariage immobilier\fg{} avec les chemins de fer, les lignes longeant les voies pour des raisons de commodité, pas de conception. En Grande-Bretagne, le problème était la densité~--- davantage de trains par mile de voie~--- et le couplage entre télégraphe et chemin de fer était intentionnel dès l'origine. \textcite{hubbard_cooke_1965} rappelle que William Fothergill Cooke avait conçu son télégraphe \emph{pour} les chemins de fer dès 1836~; en 1842, le pamphlet \emph{Telegraphic Railways} proposait formellement le \emph{block system}~--- le principe selon lequel un tronçon de voie ne pouvait être occupé que par un seul train à la fois, la vérification étant assurée par le télégraphe\footnote{Le premier \emph{block working} effectif date probablement de 1839 au \emph{Clay Cross Tunnel}, mais la datation exacte est disputée.}. Le \emph{Railway Regulation Act} de 1840 créa un organe de contrôle public, le \emph{Board of Trade}, dont les inspecteurs recommandèrent le \emph{block system} pendant quatre décennies avant qu'il ne devînt obligatoire. L'écart de quarante-sept ans entre l'invention du \emph{block system} et son adoption obligatoire en 1889~--- après la catastrophe d'Armagh, quatre-vingts morts~--- était le résultat d'une défaillance de marché au sens classique~: \dimmark{D2/D6}{+}%
le coût du télégraphe était privé, mais le bénéfice de la sécurité était en grande partie public~: les compagnies capturaient certes une part du bénéfice sous forme de réputation et d'évitement de procès, mais le risque de collision affectait d'abord les passagers, et la compagnie qui investissait dans la sécurité ne pouvait pas en facturer le prix à ses concurrentes qui ne le faisaient pas\footnote{Quinze mille instruments de télégraphe à aiguilles étaient en service sur le réseau britannique à la fin du siècle, certains jusqu'aux années~1930.}.

Ce que le cas ferroviaire faisait apparaître avec une netteté particulière, c'était la nature du problème de coordination. On peut le reformuler, avec un anachronisme assumé, comme un \emph{problème hayekien}\footnote{McCallum ne connaissait évidemment pas Hayek, dont l'article sur l'usage de la connaissance dispersée date de 1945. La reformulation est analytique~: elle éclaire la structure du problème, pas son histoire.}~: de l'information dispersée parmi des centaines d'agents répartis sur des centaines de miles devait être collectée, traitée et transformée en décisions en temps quasi réel \parencite{hayek_use_1945}. Mais le mécanisme de coordination que Hayek avait identifié comme le plus efficace~--- le système de prix~--- ne semblait pas adapté. Trois raisons structurelles l'en empêchaient~: la temporalité d'abord~--- le prix synthétise une information sur des jours ou des semaines, là où la coordination ferroviaire exigeait des décisions à la minute~; la rivalité physique absolue ensuite~--- deux trains sur la même voie signifient collision, pas compétition~; l'interdépendance séquentielle enfin~--- les agents ferroviaires sont enchaînés dans une séquence, pas en concurrence parallèle sur un marché. La réponse historique, telle que la documentent \textcite{chandler_visible_1977} et \textcite{beniger_control_1986}, ne fut pas le prix mais la hiérarchie informationnelle~--- ce que \textcite{coase_nature_1937} allait théoriser comme l'émergence de la firme là où la coordination marchande coûte plus cher que la coordination hiérarchique\footnote{La taxonomie de \textcite{hayek_use_1945}~--- la distinction entre savoir général et savoir des circonstances particulières de temps et de lieu~--- se projetait néanmoins avec exactitude sur les six formes informationnelles~: le \emph{rule book} et la standardisation codifient du savoir général~; le reporting et le télégraphe transmettent le savoir des circonstances~; la hiérarchie de McCallum est le dispositif qui articule les deux~; et la comptabilité analytique de Thomson est une tentative de construire l'équivalent du prix hayekien à l'intérieur de la firme.}. Autrement dit, le \emph{problème hayekien}~--- l'agrégation d'information dispersée~--- admettait des solutions non hayekiennes, et ces solutions étaient elles-mêmes des techniques de l'information dont le coût et la disponibilité déterminaient quelle forme organisationnelle était praticable. Le cas ferroviaire forçait ainsi la question qui allait traverser toute l'histoire des techniques de l'information~: quand le système de prix ne suffit pas à coordonner des flux matériels complexes, \dimmark{D6/D7}{*}%
quelles formes organisationnelles émergent~--- et à quel coût~?

\paragraph{Le bilan énergétique du réseau}

Or ce circuit informationnel avait un rapport au substrat énergétique que personne, à l'époque, ne pensait à mesurer~--- parce que l'asymétrie des ordres de grandeur le rendait inposable comme question. Le dispositif de \emph{process} de la division~--- les huit clercs de McCallum, les piles Daniell du télégraphe, l'éclairage du bureau~--- mobilisait au plus l'équivalent de quelques centaines de watts à quelques kilowatts de puissance dissipée en continu. En face, chaque locomotive en service développait plusieurs centaines de kilowatts, et le parc d'une division en comptait plusieurs dizaines. \dimmark{0-TER}{+}%
Quelle que soit la base de comparaison retenue, le rapport entre la puissance du circuit informationnel et la puissance du système qu'il pilotait restait inférieur de plusieurs ordres de grandeur. L'asymétrie était telle que la question du coût énergétique du circuit ne pouvait pas se poser~: l'information était, énergétiquement, gratuite.

Ce n'est pas un résultat anodin. À notre connaissance, il n'existe pas d'étude dédiée quantifiant les conséquences énergétiques des innovations informationnelles dans les chemins de fer\footnote{Résultat d'une recherche systématique sur la base Elicit (mars 2026), portant sur cinquante articles, dont dix retenus pour leur pertinence. Aucun ne traite du bilan énergétique des dispositifs informationnels ferroviaires en tant que tel.}. \textcite{field_magnetic_1992} a montré que le télégraphe contribuait à la productivité ferroviaire par la gestion serrée du capital, pas par la transmission des prix~; \textcite{fishlow_productivity_1966} a estimé que l'utilisation des capacités expliquait un sixième à un tiers de la croissance de la productivité du secteur~; \textcite{hawke_railways_1970}, cité par \textcite[p.~123--124]{miller_railway_2003}, a évalué la réduction des inventaires permise par les chemins de fer à un montant négligeable~--- entre zéro et un dixième de pour cent du revenu national en 1865. Aucun de ces travaux ne fait le dernier pas vers le bilan énergétique. Le lien entre les innovations informationnelles et la consommation d'énergie du réseau reste un trou dans la littérature~--- un trou qui s'explique par la spécialisation disciplinaire~: les historiens de la technologie ne comptent pas l'énergie, les économistes de l'énergie ne voient pas l'information, et les économistes de l'innovation traitent l'information comme analytiquement distincte de la matière.

\dimmark{D3}{+}%
\textcite{field_magnetic_1992} propose un cadre qui modifie la hiérarchie des effets. Il distingue les données de \emph{prix} --- cours du coton, actions, nouvelles politiques --- des données de \emph{quantité} --- localisation, direction et vitesse d'actifs physiques, état des stocks. Les premières conservent une valeur économique même avec un retard de quelques jours~; les secondes ne valent rien une fois que le train est arrivé ou que la viande a pourri. Le contrôle logistique en temps réel exige un dispositif à faible latence~; pour la diffusion des prix, le courrier postal eût presque suffi. Cette distinction conduit Field à réévaluer la contribution du télégraphe~: sa valeur sociale tient moins à la convergence des prix sur les bourses de matières premières qu'à l'économie de capital physique qu'il rend possible dans les systèmes de transport. Le réseau ferroviaire américain en 1890 était à 73\,\% en voie unique~; sans le télégraphe, il eût fallu doubler l'ensemble pour éviter les collisions, à un coût estimé de 957~millions de dollars --- contre 70 à 147~millions pour l'ensemble du capital télégraphique, soit un rapport de substitution de un pour sept à un pour treize. Le même raisonnement s'applique à l'industrie de la viande réfrigérée, où la périssabilité du produit interdit le stockage tampon et rend le routage en temps réel indispensable --- tandis que la sidérurgie, qui stocke cinq à trente mois de minerai, n'avait pas besoin du télégraphe pour fonctionner. La variable discriminante n'est donc pas la taille de la firme ni la complexité du réseau, mais le coût du stockage alternatif~: quand ce coût est faible, l'information optimise~; quand il est prohibitif, elle est une condition de possibilité du système.

\dimmark{D5/0-TER}{?}%
On peut en revanche conjecturer la structure du lien entre information et énergie dans ce système. Un réseau mieux coordonné consomme moins de charbon par tonne-mile~--- mais il rend aussi possible l'expansion du réseau lui-même, de trois mille miles en 1840 à plus de deux cent cinquante mille en 1900, et donc l'augmentation massive de la consommation agrégée. Si les deux effets coexistent~--- efficacité unitaire accrue et consommation totale en hausse~---, le schéma relève de ce que \textcite{jevons_coal_1865} avait identifié pour le charbon. L'hypothèse d'un \emph{effet de Jevons ferroviaire} reste une conjecture~; elle est formulée ici pour être testée au chapitre~3.

Ce cas n'est pas un septième cas de la trajectoire informationnelle~; il est le premier test du cadre analytique sur un système complet. Les dimensions que la section précédente avait identifiées isolément~--- le monopole naturel, la structure du capital, la coordination marchande, le changement technique~--- apparaissent ici en interaction. Le monopole du réseau crée un problème de coordination que le prix ne résout pas~; la coordination hiérarchique qui émerge en réponse exige un capital informationnel dont la structure diffère du capital matériel~; et le rapport de ce circuit au substrat énergétique est masqué par l'asymétrie des ordres de grandeur. Ce qui frappe, rétrospectivement, c'est la stabilité apparente du schéma~: au dix-neuvième siècle, l'énergie crée le problème~--- la vapeur dépasse la capacité de coordination~--- et l'information le résout, à un coût énergétique invisible. Le chapitre~3 examinera si, au vingt et unième siècle, ce rapport commence à s'inverser.

\emergence{D3/D7/0-TER~--- Le cas ferroviaire fournit trois résultats propres. Premièrement, l'information ne se réduit pas au télégraphe~: elle circule dans un circuit de six formes séquentielles dont le télégraphe n'est que la quatrième. Deuxièmement, il existe un seuil qualitatif au-delà duquel l'information cesse d'être un facilitateur pour devenir une condition de fonctionnement du système~--- ce que l'on pourrait nommer un \emph{seuil de McCallum}. Troisièmement, le rapport entre la puissance du circuit informationnel et celle du système coordonné est de plusieurs ordres de grandeur en faveur du second, ce qui rend le couplage analytiquement invisible. \questionch{2}{D3/0-TER~--- Cet écart se réduit-il avec les phases successives de la trajectoire~?}}

\emergence{D6/D7~--- Le cas ferroviaire force la question hayekienne~: quand le prix ne coordonne pas, quelles formes alternatives émergent~? \questionch{2}{D6/D7~--- la frontière firme-marché est-elle déterminée par les techniques d'information disponibles~?}}

\emergence{D5/0-TER~--- Le \og Jevons ferroviaire\fg{} est conjecture, pas preuve. L'information améliore l'efficacité unitaire ET permet l'expansion qui augmente la consommation agrégée. \questionch{3}{D5/0-TER~--- Le même schéma opère-t-il pour le compute AI~?}}



%%% --- BILAN §1.2 ---

\subsection*{Bilan de la section~: le télégraphe au prisme des dix dimensions}\label{subsec:bilan_telegraph}

Le tableau~\ref{tab:telegraphe-dimensions} récapitule ce que la traversée du cas télégraphe~--- réseau autonome puis cas imbriqué ferroviaire~--- a permis d'établir pour chacune des dix dimensions de la grille de lecture. Trois niveaux de certitude sont distingués~: \textbf{*}~(régularité documentée indépendamment par des sources distinctes), \textbf{+}~(mécanisme proposé, compatible avec les observations mais non testé), \textbf{?}~(piste ouverte, pas encore d'argument construit).

\begin{table}[p]
\caption{Le télégraphe au prisme des dix dimensions~: acquis, incertitudes et questions ouvertes}\label{tab:telegraphe-dimensions}
\centering\footnotesize
\renewcommand{\arraystretch}{1.15}
\begin{tabularx}{\textwidth}{@{}cXcp{4.8cm}@{}}
\toprule
\textbf{Dim.} & \textbf{Acquis du cas télégraphe} & & \textbf{Question pour la suite} \\
\midrule
D1 & Signal rival (un message occupe le circuit), information non rivale (le cours du coton est diffusable). & ? & Le numérique supprime-t-il la rivalité du canal~? \\[3pt]
D2 & Sept modèles (UK, US, FR, DE, Inde, Japon, Chine) convergent vers le monopole. Documenté indépendamment (Nonnenmacher, Kieve, Bertho, Headrick, Wu). & * & Trait des réseaux de communication, ou des infrastructures lourdes en général~? \\[3pt]
D3 & Substitution du capital linéaire (fil, poteaux) au travail (opérateurs Chappe). Quatre pays. Morse économise du cuivre au prix d'un opérateur qualifié. Le cas ferroviaire révèle deux registres~: un \emph{software} organisationnel (\emph{rule book}, hiérarchie, comptabilité) à coût faible, un \emph{hardware} matériel (télégraphe, standardisation) à coûts fixes élevés. & * & La substitution K/L de l'information est-elle du même type que celle des machines~? \\[3pt]
D4 & Opérateurs Morse~: premiers travailleurs de l'information. Huit clercs chez McCallum. Qualification non analysée. & + & L'information qualifie-t-elle ou déqualifie-t-elle~? \\[3pt]
D5 & Prix documentés (18,44~F $\rightarrow$ 1~F~; 1,55\$ $\rightarrow$ 0,40\$). Pas d'analyse d'élasticité. & ? & Effet de rebond par les tarifs~? \\[3pt]
D6 & L'information synchrone \emph{crée} le marché (Steinwender, 4~réplications). La codifiabilité filtre l'effet (Juhász). & * & La codifiabilité module-t-elle \emph{toutes} les dimensions~? \\[3pt]
D7 & Le système de prix ne coordonne pas les flux ferroviaires~: un \emph{problème hayekien} qui appelle des solutions non hayekiennes. La hiérarchie informationnelle émerge (McCallum, Chandler). Six formes séquentielles, le télégraphe n'est que la quatrième. & * & La frontière firme-marché dépend-elle des techniques d'information disponibles~? \\[3pt]
D8 & Câble transatlantique~: investissement $\rightarrow$ faillite $\rightarrow$ recapitalisation. Non formalisé. & + & Même structure que \emph{railway mania} et dot-com~? \\[3pt]
D9 & La matérialité conditionne le réseau à deux niveaux~: extraction (cuivre, gutta-percha, bois~; domination britannique) et déploiement (le réseau US trop instable ne peut faire tourner Hughes ou Baudot, qui restent cantonnés à l'Europe). & + & Même pattern pour semi-conducteurs et terres rares~? \\[3pt]
D10 & Épisodes isolés (sémaphore électrique, \emph{needle telegraph}). Pas d'argument général. & ? & Constante des réseaux de communication~? \\[3pt]
\midrule
0-TER & Asymétrie de plusieurs ordres de grandeur entre la puissance du circuit informationnel et celle du système coordonné. Matérialité extractive occultée. Aucune quantification dans la littérature. Effet de Jevons ferroviaire conjecturé. & +/? & Cet écart se réduit-il avec les phases successives de la trajectoire~? À quel seuil le couplage information-énergie devient-il visible~? \\
\bottomrule
\end{tabularx}
\renewcommand{\arraystretch}{1.0}
\end{table}

Quatre dimensions ressortent du cas télégraphe avec une évidence processuelle forte, documentée par des sources indépendantes~: la convergence vers le monopole (D2), la substitution capital-travail et la dualité \emph{software}/\emph{hardware} du capital informationnel (D3), le rôle constitutif de l'information synchrone dans la formation des marchés (D6) et l'émergence de la hiérarchie informationnelle comme réponse à la défaillance du système de prix dans le cas ferroviaire (D7). Le résultat de Steinwender sur D6 s'appuie en outre sur une stratégie d'identification économétrique répliquée dans quatre géographies. Trois dimensions sont présentes mais incomplètes~: le travail (D4), le financement (D8) et la géographie matérielle (D9)~--- les observations existent, la formalisation manque. Deux restent au stade de la piste~: la rivalité du signal (D1) et l'élasticité de la demande aux tarifs (D5). La dimension énergie (0-TER) est transversale~: elle apparaît sur trois registres~--- extraction (gutta-percha, poteaux, cuivre), invisibilité par les ordres de grandeur (ratio 1:2\,000), conjecture de Jevons ferroviaire~--- sans constituer un résultat mesurable. C'est ce silence empirique, et l'explication structurelle qu'en donne le ratio, que le chapitre~3 devra instruire.


% =============================================================================